\chapter{Solid Principles}
\label{chap:solidprinciples}

\textbf{SOLID} is a mnemonic acronym for five design principles intended 
to make object-oriented designs more understandable, flexible, and maintainable.\\
The principles are a subset of many principles promoted by 
American software engineer and instructor 
\textbf{Robert C. Martin},~\cite{martin-poO}-~\cite{martin-SOLIDstart}-~\cite{metz-sood}. \\
They were first introduced in his 2000 paper Design Principles and Design Patterns, 
while the acronym was introduced later, around 2004, 
by \textbf{Michael Feathers}~\cite{feathers-clean-architecture}. \\
Although the SOLID principles apply to any object-oriented design, 
they can also form a core philosophy for methodologies such as 
agile development or adaptive software development.

\newpage
\section{Single-Responsibility}
\label{sec:srp}
    "There should never be more than one reason for a class to change."~\cite{martin-srp}
    \newline
    In other words, every class should have only one responsibility.~\cite{martin-agileSoftDev}
    \newline
    \textbf{Smell:} a ``God'' or kitchen-sink type that mixes persistence, formatting, and notification --- any product change lands in the same file.
    \newline
    \textbf{Fix:} split along axes of change; keep orchestration thin and push each concern into its own type.
    \par\noindent\textbf{Interview tip:} say \emph{reason to change} (actor/concern), not ``small class'' --- SRP is cohesion of change, not line count.

\section{Open-Close}
\label{sec:ocp}
    "Software entities \ldots  should be open for extension, but closed for modification."~\cite{ocp}
    \newline
    New behaviour should plug in without editing stable, already-working code.
    \newline
    \textbf{Smell:} a growing \texttt{switch}/\texttt{if-else} on type tags, or scattered edits to a ``core'' module every time a variant appears.
    \newline
    \textbf{Fix:} extend via polymorphism, strategies, or plugins; seal the stable path and add new implementations beside it.
    \par\noindent\textbf{Interview tip:} name what stays closed (the hot path / published API) and what opens (a new subtype or strategy) --- OCP is not ``never touch a file again''.

\section{Liskov Substitution}
\label{sec:lsp}
    "Functions that use pointers or references to base classes must be able to use objects of derived classes without knowing it."~\cite{lsp} \\
    See also: \href{https://en.wikipedia.org/wiki/Design_by_contract}{Design by contract}. 
    \newline
    Subtypes must honour the base contract so callers can swap implementations safely.
    \newline
    \textbf{Smell:} a derived type that weakens preconditions, strengthens postconditions the wrong way, throws unexpected exceptions, or no-ops a base method (classic square/rectangle and ``not supported'' overrides).
    \newline
    \textbf{Fix:} redesign the hierarchy or split interfaces so every subtype truly is-a usable substitute; prefer composition when inheritance breaks the contract.
    \par\noindent\textbf{Interview tip:} walk one failing substitution (caller assumes base behaviour, derived surprises it) --- LSP is behavioural substitutability, not just shared method names.

\section{Interface Segregation}
\label{sec:isp}
    "Clients should not be forced to depend upon interfaces that they do not use."~\cite{isp}~\cite{martin-pap}
    \newline
    Prefer small, role-focused interfaces over fat ones that dump unused methods on every client.
    \newline
    \textbf{Smell:} a fat interface where many implementers leave methods empty, throw, or ignore half the API.
    \newline
    \textbf{Fix:} carve role interfaces (read vs write, poll vs push) so each client depends only on what it calls.
    \par\noindent\textbf{Interview tip:} point at unused methods on an implementer --- that unused surface is the ISP smell, not ``too many files''.

\section{Dependency Inversion}
\label{sec:dip}
    "Depend upon abstractions, [not] concretions."~\cite{dip}~\cite{martin-pap}
    \newline
    High-level policy should not hard-wire low-level details; both should depend on shared abstractions.
    \newline
    \textbf{Smell:} domain logic that constructs concrete I/O, DB, or vendor SDKs directly --- swapping a dependency means editing the core.
    \newline
    \textbf{Fix:} introduce an interface (or port) owned by the high level; inject the concrete adapter at the boundary.
    \par\noindent\textbf{Interview tip:} draw who owns the abstraction --- DIP inverts the usual ``UI depends on DB'' arrow so details plug into policy, not the reverse.
